\section{Можно ли устранить отставание схемой или протоколом}
\label{sec:protocol}
Полезность недавней истории ещё не означает, что исходные взаимодействия принципиально необходимы. Альтернативное объяснение состоит в неудачном способе записи состояния. Отдельный инструмент обновления отличается от совместного контракта SKILL.state, а общая схема может плохо представлять работу в терминале. Для проверки этих объяснений внутри гибридной системы сравниваются сочетания tool/inline-протоколов и исходной/адаптированной схем (табл.~\ref{tab:protocol}).

\begin{table}[htbp]
\centering\small
\caption{Протокол и схема внутри гибридного состояния без недавней истории,
в тех же двух панелях, что и табл.~\ref{tab:composition}. Гибридные сводки
не позволяют восстановить парные отношения стоимости, пиковый промпт
и cache hit (--).}
\label{tab:protocol}
\resizebox{\linewidth}{!}{\input{../paper/tables/new-protocol}}

\medskip
\resizebox{\linewidth}{!}{\input{../paper/tables/new-protocol-resources}}
\end{table}

Адаптированная схема при отдельном инструменте даёт \TaskAdaptedRate\% против \HybridRate\% у опорного варианта. Наблюдаемое улучшение небольшое, а медианная стоимость возрастает с \$\HybridCostMedian{} до \$\TaskAdaptedCostMedian{}. Для восстановления контрольной успешности этого изменения недостаточно.

Inline-протокол с исходной схемой даёт \InlineRate\% и медианную стоимость \$\InlineCostMedian{}. Совместное изменение протокола и схемы даёт \InlineAdaptedRate\% при стоимости \$\InlineAdaptedCostMedian{}. Вместо ожидаемого устранения двух ограничений наблюдается ухудшение. Описательная разность эффектов inline-протокола между схемами составляет $\ProtocolSchemaInteraction$ процентного пункта. Проверка такого взаимодействия с учётом задач требует совместных исходов, которых нет в доступной агрегированной сводке.

У inline-вариантов больше медианный выход. Совместное формирование обновления и действия поэтому не оказалось бесплатной заменой отдельного инструмента. Расход может включать текст обновлений, исправления формата и другие изменения траектории; доступные агрегаты не разделяют эти составляющие.

Это сравнение исключает простое объяснение, согласно которому достаточно применить именно проверенную схему или совместный формат ответа, чтобы восстановить результат полной истории. Оно не исключает более удачных схем и реализаций. Измеренные изменения отличаются от возврата недавней истории: последний повышает успешность внутри обеих систем памяти, тогда как выбранные изменения представления не дают сопоставимого восстановления.
